%%%PAGE 175 rot=0 legibility=good summary=测动摩擦因数μ的整体法与图像法、斜面测μ不用天平、验证机械能守恒（图像法、光电门）
%%%BLOCK id=175-01 ch=03 sec=测动摩擦因数·实验 kind=method exp=1 alt=none cont=none y=0.12-0.29
\textbf{(四)\ 测 $\mu$}

%FIG p175_a 0.05 0.17 0.26 0.30
\notefig[0.25]{figures/p175_a.png}

整 $mg-\mu Mg=(M+m)a$ % 行首「整」

法一\quad $\mu=\dfrac{mg-(M+m)a}{Mg}$

%FIG p175_b 0.60 0.185 0.70 0.247
\notefig[0.25]{figures/p175_b.png}

测 $k$、斜截距 \struck{\unclear{…}} % CHECK: 「测」与「k」之间有一竖线（可能是分隔符）；行末有一处涂黑的字

或用数据代入求平均
%%%END
%%%BLOCK id=175-02 ch=03 sec=测动摩擦因数·实验 kind=method exp=1 alt=none cont=none y=0.29-0.55
%FIG p175_c 0.04 0.29 0.31 0.47
\notefig[0.27]{figures/p175_c.png}

令 $m'+m=m_0$

\begin{align*}
& mg-\mu(M+m')g=(M+m'+m)a \\
& mg-\mu(M+m_0-m')g=(M+m_0)a % CHECK: 手写为 m'（带撇），按 m'+m=m_0 此处应为 m
\end{align*}

\[
a=\dfrac{(1+\mu)g}{M+m_0}\,m-\mu g % CHECK: 分母手写像小写 m+m_0，按上文应为 M+m_0
\]

%FIG p175_d 0.60 0.37 0.985 0.55
\notefig[0.4]{figures/p175_d.png}
%%%END
%%%BLOCK id=175-03 ch=03 sec=测动摩擦因数·实验 kind=method exp=1 alt=none cont=none y=0.56-0.72
哪些实验要天平（能约就不用天平，不能约的用天平）\qquad 先写式子再看函数

%FIG p175_e 0.165 0.605 0.335 0.705
\notefig[0.25]{figures/p175_e.png}

\begin{align*}
& mg\,\frac{h}{L}-\mu mg\,\frac{\sqrt{L^2-h^2}}{L}=ma \\
& \mu=\dfrac{gh-aL}{g\sqrt{L^2-h^2}}
\end{align*}
%%%END
%%%BLOCK id=175-04 ch=06 sec=验证机械能守恒·实验 kind=method exp=1 alt=none cont=next y=0.73-0.97
验机械能守恒；能量 $E_p=\frac{1}{2}kx^2$

\hspace{2em}本质—式子函数

\hspace{2em}验证$\to$图像$\to$直线

$mgh=\frac{1}{2}mv^2$

%FIG p175_f 0.53 0.725 0.71 0.83
\notefig[0.25]{figures/p175_f.png}

\begin{center}
\begin{tabular}{|c|p{6cm}|}
\hline
$h$ & \\
\hline
$v$ & \\
\hline
$v^2$ & \\
\hline
\end{tabular}
\end{center}

\textbf{装置1}

%FIG p175_g 0.21 0.855 0.39 0.93
\notefig[0.25]{figures/p175_g.png}

$mgh=\frac{1}{2}mv^2$\quad 不用天平

$\downarrow$ 光电门 $v=\dfrac{d}{\Delta t}$

%FIG p175_h 0.73 0.825 0.98 0.945
\notefig[0.25]{figures/p175_h.png}
%%%END
%%%PAGEEND 175
